\section{Introduction}\label{sec:intro}

A projective Calabi--Yau threefold can have only locally smoothable
singularities and nevertheless admit no smoothing.  For ordinary double
points, the obstruction is expressed by the exceptional curves of a small
resolution: a smoothing exists precisely when their classes satisfy a
relation
\begin{equation}\label{eq:friedman}
 \sum_j\lambda_j[C_j]=0,\qquad \lambda_j\ne0\quad\text{for every }j.
\end{equation}
This is Friedman's criterion, together with the unobstructedness and
necessity results discussed in \cite{Friedman86,Namikawa94,RT09};
\cite[Thm.~1.3]{RT09} states the smoothing criterion itself.
We prove an analogous criterion when nodes occur together with
anticanonical cones over smooth del Pezzo surfaces of degrees five, six and seven,
and over $\PP^1\times\PP^1$, of degree eight.
The degree-six cone has a two-dimensional smoothing component, and a
nonzero vector on that component need not be a smoothing direction.
Three separate nonvanishing conditions are necessary.  An explicit
nonsmoothable threefold shows that they cannot be replaced by a single
nonzero-vector condition, even if the local components may be chosen anew.

\subsection{The smoothing criterion}
Let $X$ be a connected normal projective complex threefold with
$\omega_X\cong\cO_X$ and $H^1(X,\cO_X)=0$.  Suppose every singularity is
an ordinary double point or is analytically isomorphic to the
anticanonical cone over a smooth del Pezzo surface of degree five, six or seven,
or over $\PP^1\times\PP^1$.
The assumption concerns these exact analytic germs.
Let $Y_0\to X$ be the smooth crepant analytic resolution obtained from
chosen small resolutions of the nodes and the standard divisorial
resolutions of the cones.  It is compact and Moishezon.

A \emph{branch profile} $S$ chooses a reduced local deformation component
at each singular point.  At a node or a $\PP^1\times\PP^1$ cone the base
is a line; at a degree-five point it is a smooth four-dimensional germ;
at a degree-seven
point there is one reduced line; at a degree-six point there are a line
and a plane, denoted $A_1$ and $A_2$.  These names agree with the root
subspaces in the del Pezzo lattice.  For the link $L_p$ of a cone point,
the circle-bundle projection identifies
\[
 H_2(L_p;\QQ)\cong K_{E_p}^{\perp}\subset H_2(E_p;\QQ).
\]
Write $R_{p,S_p}$ for the nodal line or the selected root subspace, and set
\begin{equation}\label{eq:relationkernel}
 K_S=\ker\left\{\bigoplus_pR_{p,S_p}
                         \longrightarrow H_2(Y_0;\QQ)\right\}.
\end{equation}
At a $\PP^1\times\PP^1$ cone the selected subspace is the entire
link-homology line $\QQ(f_1-f_2)$, where $f_1,f_2$ are the ruling
classes on the exceptional surface.
At a degree-six point, mark $E_p=\mathrm{Bl}_3\PP^2$ by the line class
$h$ and exceptional classes $e_1,e_2,e_3$.  Its $A_2$ subspace is
\[
R_{p,A_2}=\left\{\mu_1e_1+\mu_2e_2+\mu_3e_3:
                                  \mu_1+\mu_2+\mu_3=0\right\}.
\]
At a degree-five point mark $E_p=\operatorname{Bl}_4\PP^2$
by $h,e_1,\ldots,e_4$.  The selected subspace is the entire
$K_{E_p}^\perp$, and its five conic-pencil classes are
\[
 f_i=h-e_i\ (1\le i\le4),\qquad
 f_5=2h-e_1-e_2-e_3-e_4.
\]
Put $\lambda_i(\kappa_p)=\kappa_p\cdot f_i$.  Their sum is
zero because $\sum_i f_i=-2K_{E_p}$.

\begin{thmletter}[smoothing criterion]\label{thm:criterionintro}
The threefold $X$ has a convergent one-parameter smoothing inducing $S$
if and only if $K_S$ contains a vector nonzero at every rank-one summand
and satisfying $\mu_1\mu_2\mu_3\ne0$ at every $A_2$ summand
and $\prod_{i=1}^5\lambda_i\ne0$ at every degree-five point.
\end{thmletter}

Equivalently, none of the stated scalar functionals vanishes identically
on $K_S$.  Thus the criterion is a finite linear-algebra test for each
of the finitely many profiles.  When neither degree-five nor degree-six
points occur, its form is the same as the nodal criterion.  At an $A_2$
point the three excluded lines are precisely the local discriminant.
Necessity holds for arbitrary smoothing arcs, including arcs tangent to
that discriminant to higher order.
At degree five, the five excluded hyperplanes are likewise the
exact local discriminant, and necessity holds for arbitrary arcs.

The integrability statement is stronger than the existence assertion.
Fix analytic local classifying maps for a semiuniversal deformation of
$X$, and let $\Def^S(X)$ be the scheme-theoretic inverse image of the
chosen local branches.  Denote the tangent sheaf by $\Theta_X$.

\begin{thmletter}[smooth selected deformation spaces]\label{thm:smoothintro}
If $K_S\to R_{p,S_p}$ is nonzero at every degree-six point, then
$\Def^S(X)$ is smooth, of dimension
\[
 h^1(X,\Theta_X)+\dim_{\QQ}K_S.
\]
Every tangent vector to this space is the tangent of a convergent
one-parameter deformation on the selected branches.
\end{thmletter}

No nonzero-projection hypothesis is required at nodes, degree-five or degree-seven
points, or $\PP^1\times\PP^1$ cones in this theorem.  The hypothesis
is automatic when the smoothing
criterion holds.  Smoothness of a selected base does not, by itself,
assert the existence of a smooth fibre.

\subsection{Scope among del Pezzo cones}\label{sec01:degree-scope}

The degree-seven and degree-six cases arise from the pentagonal and
hexagonal faces in the toric examples.  Both associated cones are smoothable and are
not complete intersections.  Their local deformation spaces exhibit
the two phenomena treated here: distinct smoothing components in
degree six, and a nonreduced structure transverse to the smoothing
line in degree seven.  Theorems~\ref{thm:criterionintro}
and~\ref{thm:smoothintro} also allow the degree-five cone and the
anticanonical cone over $\PP^1\times\PP^1$, of degree eight.
These hypotheses specify the
local comparisons established below; smoothable del Pezzo cones also
occur in the other degrees.

For a smooth del Pezzo surface $E$, the anticanonical cone means
\[
 \operatorname{Spec}\bigoplus_{m\geq0}H^0(E,-mK_E).
\]
This convention includes the weighted
anticanonical models in degrees one and two.  The classical local
facts in Table~\ref{tab01:degree-scope} are recalled in
\cite[\S\S3.3, 5.5 and Example~2.19]{Gross97rev}; the branch structures
in degrees six and seven follow from \cite{Altmann97}.
Proposition~\ref{prop03:degree5family} proves the full local-base
and discriminant statements in degree five.

% Canonical degree-scope table.  Keep its final column synchronized with
% the theorem statements and the degree-five/degree-eight research TODO.
% Local smoothability and an established global criterion are distinct.
\begin{table}[htbp]
\centering
\caption{Anticanonical cones over smooth del Pezzo surfaces and the
scope of the present results.  The local assertions alone do not imply
smoothability of a compact threefold containing these germs.}
\label{tab01:degree-scope}
\small
\renewcommand{\arraystretch}{1.14}
\begin{adjustbox}{max width=\textwidth, center}
\begin{tabular}{@{}>{\raggedright\arraybackslash}p{0.17\textwidth}
                  >{\raggedright\arraybackslash}p{0.41\textwidth}
                  >{\raggedright\arraybackslash}p{0.32\textwidth}@{}}
\toprule
Degree / surface & Local deformation behavior & Role in the present results \\
\midrule
$1$--$3$
  & Hypersurfaces; smoothable, with smooth local bases.
  & Classical complete-intersection cases; outside the stated mixed criterion. \\
\addlinespace
$4$
  & Intersection of two quadrics; smoothable, with smooth local base.
  & As in degrees one to three. \\
\addlinespace
$5$
  & Smooth four-dimensional base; five discriminant hyperplanes.
  & Covered by Theorems~\ref{thm:criterionintro}
    and~\ref{thm:smoothintro}; five nonzero conic-pencil pairings. \\
\addlinespace
$6$
  & A smoothing line and a smoothing plane; three discriminant lines in the plane.
  & Both components covered by Theorems~\ref{thm:criterionintro}
    and~\ref{thm:smoothintro}. \\
\addlinespace
$7$
  & One reduced smoothing line, with an obstructed transverse tangent.
  & Covered by Theorems~\ref{thm:criterionintro}
    and~\ref{thm:smoothintro}. \\
\addlinespace
$8$, $\PP^1\!\times\!\PP^1$
  & Smooth one-dimensional local base; noncentral fibres smooth.
  & Covered by Theorems~\ref{thm:criterionintro}
    and~\ref{thm:smoothintro}; coefficient on the ruling difference. \\
\addlinespace
$8$, $\mathbb F_1$
  & No smoothing; local base $\operatorname{Spec}\CC[t]/(t^2)$.
  & Excluded by local nonsmoothability. \\
\addlinespace
$9$, $\PP^2$
  & Rigid.
  & Excluded by local nonsmoothability. \\
\bottomrule
\end{tabular}
\end{adjustbox}
\end{table}

The degree-eight cone over $\PP^1\times\PP^1$ already occurs as the
partial resolution of a degree-seven point.
Lemma~\ref{lem03:quadric} proves smoothness of its full local base and
identifies its tangent line with the difference of the ruling classes.
When such points occur on the original threefold, the partial
modification is the identity near them.  Their Milnor fibres retract
to $\mathbb{RP}^3$; the descended nodal period supplies the necessity
argument in Lemma~\ref{lem04:periodnecessity}.  Thus they require one
nonzero coefficient on the ruling difference, with no additional
hypothesis in the selected-base smoothness theorem.

The degree-five surface is not toric.  Totaro's Bott vanishing
theorem \cite{totaro2020} supplies its local Hodge comparison.
The four-parameter family of Grassmannian sections and its
boundary periods identify the five conic-pencil conditions and
prove necessity at every order.  Degree-five points also remain
unchanged on the partial model.  Thus the present criterion covers
all locally smoothable anticanonical cones over smooth del Pezzo
surfaces that are not complete intersections.
For degrees one to four, the existing complete-intersection theory
does not by itself give the stated criterion for mixtures with the
non-complete-intersection cones considered in this paper.


\Needspace{18\baselineskip}
\subsection{A nonsmoothable example and its deformation space}
The distinction between a nonzero local vector and a smoothing direction
cannot be removed by changing the branch profile.

\begin{samepage}
\begin{thmletter}\label{thm:counterintro}
There is a projective Calabi--Yau threefold with thirty nodes and two
degree-six cone points such that a relation in $K_S$ is nonzero at every
singular point for either mixed profile, but the threefold has no
smoothing on any profile.  Its semiuniversal analytic local ring is
\[
 \CC\{z_1,\ldots,z_{32},x,y,u,v,b\}/(xy,xb,uv,ub).
\]
It is reduced, has tangent dimension $37$, and has four smooth
irreducible components of dimensions $34,34,34,35$.  All their
scheme-theoretic intersections are smooth.
\end{thmletter}
\end{samepage}

The example is a general anticanonical hypersurface in the toric
fourfold associated with an explicit reflexive polytope with $25$
vertices.  Three divisor identities prove nonsmoothability independently
of the sufficiency argument.  They force an old nodal coefficient to
vanish on each uniform profile and one of the three $A_2$ coefficients
to vanish on each mixed profile.  The full deformation ring then follows
from Theorem~\ref{thm:smoothintro}, the local branch equations, and a
convergent change of coordinates.  A prescribed tangent integrates
precisely when the four displayed quadratic expressions vanish.

\subsection{The geometric argument}
The first-order comparison identifies the image of the global tangent
space in the sum of the local tangent spaces with the kernel of the
link-homology map.  Contraction with one global volume form, local
cohomology with supports, and the $\partial\bar\partial$-lemma give this
comparison for the full local tangent spaces, including the transverse
infinitesimal direction at a degree-seven cone.

To integrate the selected directions, we replace the cone points by
proper crepant partial resolutions.  The degree-seven cone is replaced
by the anticanonical cone over $\PP^1\times\PP^1$; the $A_1$ and $A_2$
choices at a degree-six cone are replaced by two and three nodes,
respectively.  Degree-five and original $\PP^1\times\PP^1$
cones remain unchanged.
Each resulting singularity has a smooth full local
deformation base.  Friedman's generalized $T^1$ lifting theorem
\cite[Thm.~3.10]{friedman2026} therefore proves unobstructedness of the
partial model.  An analytic blow-down lemma and a lift of locally
trivial global tangents transfer this smoothness to the selected
deformation space of $X$.  The partial model need not be projective;
compact Moishezon resolutions supply the required Hodge theory.

For necessity at an $A_2$ point, an explicit simultaneous partial
resolution of the rank-one matrix family is an isomorphism over every
noncentral fibre.  Pulling it back along an arbitrary smoothing arc
replaces the central cone by three nodes without changing the smooth
fibres.  Nodal necessity then supplies three nonzero coefficients;
intersection with the exceptional del Pezzo surface makes their sum
zero.  This proves the stronger $A_2$ condition at every order.
At a $\PP^1\times\PP^1$ cone, the vanishing cycle is the quotient
of the nodal sphere by the antipodal involution; its nonzero period
gives the corresponding rank-one necessity condition.
At degree five, the four-parameter Grassmannian section family
has five discriminant hyperplanes.  Its boundary carries every
absolute Milnor cycle.  The resulting holomorphic boundary periods
identify the five conditions and prove necessity along arbitrary arcs.

\subsection{Toric constructions and their limits}
The general theorems also distinguish abstract deformations from
deformations induced by a toric ambient space.  For Batyrev
hypersurfaces \cite{Batyrev94}, we study the ambient families obtained
from an admissible Minkowski decomposition at one primitive lattice
degree, using Altmann's local theory and the global construction of
Ilten--Vollmert \cite{Altmann97,IV12}.  Definition~\ref{def:singledegree}
specifies these families.

The hypersurface $X_{19}$ has fourteen nodes and one degree-seven cone
point.  Corollary~\ref{cor:x19smooth} proves that its full deformation
space is smooth of dimension $30$ and that it is smoothable.
Nevertheless, Theorem~\ref{thm:A} excludes a smoothing by every
single-degree ambient family through the relative anticanonical system.
Its proof combines Altmann's grading calculation with a propagation
rule on the edges of the two facets containing the pentagonal face.
Definition~\ref{def:locked} calls a facet \emph{locked} when this rule
forces all its edge-dilation parameters to agree.  The obstruction is
uniform over the entire degree lattice.

The complementary construction is $X_9$, with two nodes and one
degree-seven cone point.  Theorem~\ref{thm:B} gives an explicit ambient
smoothing at degree $(0,1,-1,-1)$.  Theorem~\ref{thm:C} classifies all
$77$ admissible polytopes with at most nine vertices and a pentagonal
face: $76$ give nonsmoothable hypersurfaces, and the remaining one is
$X_9$.  The larger edge-propagation census describes the frequency of
the ambient obstruction, without asserting that its separate
degree-coverage condition holds throughout that census.

Sections~\ref{sec:deformation}--\ref{sec04:smoothing} prove the general
theorems, and Section~\ref{sec05:counterexample} gives the counterexample
and its full deformation ring.  Sections~\ref{sec:germs}--\ref{sec:thmB}
develop the toric applications.  Appendix~\ref{app:cert} records the
exact computations and the data from which they can be reproduced.
The results leave open smoothness of selected bases whose degree-six
projection vanishes, and the structure of unrestricted deformation
bases in general, including possible degree-seven nilpotents.
